\chapter{API Documentation}
\label{appendix:api-docs}

\section{REST API Endpoints}
\label{app:rest_api}

The {\name} system provides a comprehensive REST API for programmatic access to all functionality. This section documents all available endpoints, request/response formats, and usage examples.

\subsection{Base Configuration}

\begin{itemize}
\item \textbf{Base URL:} \texttt{http://localhost:5000}
\item \textbf{Content-Type:} \texttt{application/json}
\item \textbf{Authentication:} None (local deployment)
\item \textbf{Rate Limiting:} 100 requests per minute
\end{itemize}

\section{Core Processing Endpoints}
\label{app:core_endpoints}

\subsection{Process URLs}

\textbf{Endpoint:} \texttt{POST /process-urls}

Primary endpoint for processing website URLs and generating clustering analysis.

\textbf{Request Parameters:}

\begin{table}[h]
\centering
\caption{Process URLs Request Parameters}
\begin{tabular}{|l|l|l|l|}
\hline
\textbf{Parameter} & \textbf{Type} & \textbf{Required} & \textbf{Description} \\
\hline
\texttt{urls} & Array[String] & Yes & List of URLs to process \\
\texttt{levels} & Array[Integer] & Yes & DOM analysis levels [1,2,3] \\
\texttt{mode} & Integer & No & Processing mode (default: 0) \\
\texttt{custom\_ttl\_hours} & Integer & No & Cache TTL in hours \\
\hline
\end{tabular}
\end{table}

\textbf{Example Request:}
\begin{lstlisting}[language=json]
{
  "urls": [
    "https://example.com",
    "https://test.com",
    "https://sample.org"
  ],
  "levels": [1, 2, 3],
  "mode": 0,
  "custom_ttl_hours": 24
}
\end{lstlisting}

\textbf{Response Format:}
\begin{lstlisting}[language=json]
{
  "status": "success",
  "data": {
    "processing_time": 15.7,
    "cache_hit_rate": 73.2,
    "total_elements_processed": 1247,
    "clusters_generated": 12,
    "outliers_count": 89,
    "dbcv_scores": {
      "level_1": 0.347,
      "level_2": 0.392,
      "level_3": 0.421
    },
    "visualization_path": "/results/viz_1697834567890.html"
  }
}
\end{lstlisting}

\subsection{Site Similarity Analysis}

\textbf{Endpoint:} \texttt{POST /site-similarity}

Computes cosine similarity between website structures.

\textbf{Request Example:}
\begin{lstlisting}[language=json]
{
  "urls": ["https://site1.com", "https://site2.com"],
  "levels": [1, 2],
  "similarity_metric": "cosine"
}
\end{lstlisting}

\textbf{Response Format:}
\begin{lstlisting}[language=json]
{
  "status": "success",
  "data": {
    "similarity_matrix": [
      [1.0, 0.847],
      [0.847, 1.0]
    ],
    "similarity_score": 0.847,
    "analysis_details": {
      "common_features": 156,
      "total_features": 184,
      "feature_overlap_percentage": 84.7
    }
  }
}
\end{lstlisting}

\section{Cache Management Endpoints}
\label{app:cache_endpoints}

\subsection{Cache Statistics}

\textbf{Endpoint:} \texttt{GET /cache/stats}

Returns comprehensive cache performance metrics.

\textbf{Response Format:}
\begin{lstlisting}[language=json]
{
  "status": "success",
  "data": {
    "cache_performance": {
      "overall_hit_rate": 85.3,
      "complete_cache_hit_rate": 73.2,
      "partial_cache_hit_rate": 42.1,
      "total_requests": 2847,
      "cache_hits": 2428,
      "cache_misses": 419
    },
    "storage_statistics": {
      "total_cache_size_mb": 234.7,
      "complete_cache_size_mb": 178.3,
      "partial_cache_size_mb": 56.4,
      "cache_entries": 847
    },
    "performance_metrics": {
      "average_response_time_cached": 0.15,
      "average_response_time_uncached": 8.7,
      "processing_time_reduction_percent": 71.7,
      "network_requests_saved_percent": 70.1
    }
  }
}
\end{lstlisting}

\subsection{Cache Operations}

\textbf{Clear Cache:} \texttt{DELETE /cache/clear}

\textbf{Cache Cleanup:} \texttt{POST /cache/cleanup}

\textbf{Cache Configuration:} \texttt{GET /cache/config}

\section{Error Handling}
\label{app:error_handling}

All API endpoints follow a consistent error response format:

\textbf{Error Response Format:}
\begin{lstlisting}[language=json]
{
  "status": "error",
  "error": {
    "code": "INVALID_PARAMETERS",
    "message": "The provided URLs list is empty",
    "details": {
      "parameter": "urls",
      "expected": "non-empty array",
      "received": "empty array"
    },
    "timestamp": "2024-12-16T10:30:45Z"
  }
}
\end{lstlisting}

\subsection{Common Error Codes}

\begin{table}[h]
\centering
\caption{API Error Codes}
\begin{tabular}{|l|l|l|}
\hline
\textbf{Code} & \textbf{HTTP Status} & \textbf{Description} \\
\hline
\texttt{INVALID\_PARAMETERS} & 400 & Invalid request parameters \\
\texttt{URL\_FETCH\_ERROR} & 422 & Failed to fetch website content \\
\texttt{PROCESSING\_ERROR} & 500 & Internal processing error \\
\texttt{CACHE\_ERROR} & 500 & Cache system error \\
\texttt{CLUSTERING\_ERROR} & 500 & Clustering algorithm error \\
\texttt{TIMEOUT\_ERROR} & 504 & Request timeout exceeded \\
\hline
\end{tabular}
\end{table}

\section{API Client Examples}
\label{app:client_examples}

\subsection{Python Client Example}

\begin{lstlisting}[language=Python]
import requests
import json

def process_urls(urls, levels=[1, 2, 3]):
    endpoint = "http://localhost:5000/process-urls"
    payload = {
        "urls": urls,
        "levels": levels,
        "mode": 0
    }
    
    response = requests.post(endpoint, json=payload)
    
    if response.status_code == 200:
        result = response.json()
        print(f"Processing completed in {result['data']['processing_time']}s")
        return result
    else:
        print(f"Error: {response.status_code}")
        return None

# Usage example
urls = ["https://example.com", "https://test.com"]
result = process_urls(urls)
\end{lstlisting}

\subsection{JavaScript Client Example}

\begin{lstlisting}[language=JavaScript]
async function processUrls(urls, levels = [1, 2, 3]) {
    const endpoint = 'http://localhost:5000/process-urls';
    const payload = {
        urls: urls,
        levels: levels,
        mode: 0
    };
    
    try {
        const response = await fetch(endpoint, {
            method: 'POST',
            headers: {
                'Content-Type': 'application/json',
            },
            body: JSON.stringify(payload)
        });
        
        if (response.ok) {
            const result = await response.json();
            console.log(`Processing completed in ${result.data.processing_time}s`);
            return result;
        } else {
            throw new Error(`HTTP error! status: ${response.status}`);
        }
    } catch (error) {
        console.error('Error:', error);
        return null;
    }
}

// Usage example
const urls = ['https://example.com', 'https://test.com'];
processUrls(urls).then(result => {
    if (result) {
        console.log('Success:', result);
    }
});
\end{lstlisting} 